\section{Intro}

\subsection{Background and Motivation}

"There is no denying the fact that electricity is crucial for the world economy to thrive in this and the coming centuries. The demand for electricity increases not only because the human population increases, but also because the social-economic activities of humans are rapidly shifting from manual processes to automated processes, which are essentially driven by electricity. Therefore, electricity becomes inevitable for the sustainability of modern civilization because it has found its way to the root of many human establishments [1,2]. Electricity can be generated from various sources and its generation from many sources is desired to boost the total generation capacity to meet the ever-growing demand. As electrical power cannot be stored, instantaneous balancing becomes a necessity and it is considered a key design parameter in power systems [3]." \cite{salman_review_2022}

"The evolution of the power system has significant impact on the operation and planning of the future power system. Five major global trends influencing the evolution of the power system are [1]: Decarbonisation - Decreasing the carbon footprint from electric power production. Decentralisation - Transition from few and large, centralized, power plants to many smaller, decentralised, power production units. Integration – Increasingly integrated electricity markets, greater interconnection of previously independent grids, and more integrated energy systems including sector coupling. \cite{emil_hillberg_flexibility_2019}.


\subsection{Problem Statement}
"The energy transition leads to a shift from conventional, controllable power plants to renewable, variable generation resources such as wind and solar power plants. According to [1,2], the existing renewable power capacity worldwide (including hydro power) nearly tripled from 1150 GW to 3146 GW between 2008 and 2021. The variability and uncertainty of renewable generation leads to an increased need for flexibility resources." \cite{lechl_review_2023}




\section{Theory}
\subsection{Demand side management}
In [8], DSM is described as activities to activate the demand side, comprising actions such as energy efficiency, savings, self-production and load management. Further, load management techniques and DR are examples of DSM solutions. According to [38], there are six typical versions of DSR, such as conversion and energy efficiency, load shifting, peak clipping, valley filling, flexible load shape and electrification. \cite{degefa_comprehensive_2021}


\section{Discussion}
\subsection{Challenges and Barriers to Implementing Flexibility Solutions}
